\section{LLM LOOP}
\label{app:llm-loop}

This appendix records a separate, executable whole-solver recovery experiment.
The question is whether a model given a problem class, ordinary algorithm
references, and a set of legal moves can assemble one of the complete methods
used in this paper. The target here is the hyperspectral unmixing (HSI) solver,
not a new speed record. All candidate programs, model requests, failure records,
source hashes, manifests, and independent certificate outcomes are retained
under \path{results/executable_discovery_20260926/}. Unless stated otherwise,
the target is a relative Frank--Wolfe gap of $10^{-4}$ \emph{for every column}
and simplex feasibility error at most $10^{-9}$.
Execution feedback inside an LLM evolutionary loop is prior
art~\citep{ye2024reevo}; the controlled question here is whether the withheld
expert hints and the executable move space permit verified method recovery at
a fixed fresh-evaluation budget.

\subsection{Problem, moves, and judgment}
For a shared $D\in\mathbb R^{n\times p}$ and columns $y_j$, the model must
return the complete matrix $X$ for
\begin{equation}
  \min_{x_j\geq0,\,{\bf1}^\top x_j=1}
       \tfrac12\lVert D x_j-y_j\rVert_2^2,\qquad j=1,\ldots,m.
\end{equation}
The real spectral dictionary has shape $188\times498$. The fixed judge uses
host float64 arithmetic on the original, unmodified $D,Y$: it checks shape,
finiteness, simplex feasibility, and
$[\langle\nabla f_j(x_j),x_j\rangle-\min_i(\nabla f_j(x_j))_i]/
 [1+|f_j(x_j)|]\leq10^{-4}$ for every $j$.
A failed or over-budget case receives the full time cap in the search score.
This independent output certificate cannot be weakened by a generated stopping
rule. A candidate first encounters three numerical admission cases (identity,
rank-deficient, and zero design). Admission, setup, JIT compilation, transfers,
iteration, output reconstruction, and repeated host certificates are all counted
in its algorithm clock; JAX runtime initialization is separately recorded as
process overhead. Every timed evaluation uses a fresh worker process with its
persistent compilation cache disabled.

The interface has eight executable layers: (1) quadratic formulation or
forward gradient; (2) FBS, DRS, or direct equality-split ADMM update;
(3) carried state and inverse, Cholesky, recomputed, or host-computed factor;
(4) optional retirement of independently accepted columns with full-batch
reconstruction; (5) optional finite FBS-to-DRS handoff; (6) internal float32 or
float64; (7) CPU or GPU execution; and (8) compiled chunk length and unrolling.
A separate certificate-schedule choice requests a host check every chunk or
uses a cheap device screen between host checks. Step size, resolvent parameter,
and relaxation are also selectable. The compiler supplies the named standard
operators; the proposer chooses a composition and parameters from this finite
space. Thus recovery is \emph{automatic assembly from supplied moves}, not the
invention of ADMM or of a new simplex projection. These experiments check the
base FBS/DRS constructions and the numerical outputs, but the full direct ADMM,
retirement, and certificate-scheduling adapter does not carry a new
whole-program convergence proof. This experiment must not be conflated with the
paper's formal admissibility claims for its original certified genome calculus.

\subsection{Data, hardware, and model access}
The final recovery search uses disjoint training pixels from the same Cuprite
image: $m=4{,}096$ and $16{,}384$ with the shared $188\times498$ dictionary,
plus a synthetic $96\times160$ design with $32$ right-hand sides. The search
budget is $120$\,s per case. A sealed test uses \emph{different Cuprite pixels}
in batches of $6{,}144$ and $10{,}000$, plus a synthetic $192\times320$
design with $128$ right-hand sides and stronger conditioning; each is run three
times at the same $120$\,s cap. Pixel indices and input hashes were frozen in
the manifests before this search. Real-pixel holdouts share a spectral library;
this is within-library transfer, alongside a synthetic dimension and
conditioning shift, not an unseen real dataset.

There are six A100 80\,GB devices, four on node~0 and two on node~1. The
separately hosted H100 80\,GB serves the frozen Qwen 27B proposer and does not
also run timed candidates. The local RTX~4090 carries an existing smaller-model
service, not this assay. Each candidate chain owns an A100 slot; per-GPU file
locks prevent timed jobs from overlapping. The model sees only the problem,
training case dimensions, its own earlier configurations and measured outcomes,
the action grammar, and standard textbook equations and executable NumPy
examples of projected gradient, ADMM, DRS, quadratic proximal steps, and simplex
projection. It never receives the manually authored completed HSI solver or
sealed inputs. The \texttt{full\_cards} arm additionally sees four short expert
hints: split the quadratic from the simplex constraint, reuse a factor across
right-hand sides, retire accepted columns while retaining all outputs, and
price compiled grouping and certificate cadence. The \texttt{withheld\_cards}
arm sees the same standard references and action support without these hints.
The \texttt{resample} arm does not receive execution history; the random arm
samples the same canonical parameter space. All four arms completed two
prescribed seeds and eight fresh executable candidates per seed. Duplicates and malformed
responses are logged separately, with a hard proposal-call cap. No model
weights are updated in this simple loop.

\subsection{How the result was reached}
The original free-source experiment demonstrated why mere prompting was
insufficient. In a $96$-candidate whole-module pilot (four conditions, two
seeds, twelve fresh programs each), only three generated programs solved both
small training cases: one with no standard references and two from independent
resampling. The card-assisted and withheld-card arms solved none. The initial
pilot also had a seed-baseline output-path collision, so its candidate outcomes
are useful failure records but its baseline timings are not used for a
performance claim. A later function-patch experiment produced an accelerated
projected-gradient candidate and, in one withheld-card lineage, an LLM-written
CPU active-set solver. Seven valid descendants in that lineage are variations
of the same family, not seven independent discoveries. The active-set program
certified five of five Cuprite-512 validation repetitions, but failed all five
Cuprite-8192 repetitions under the $60$\,s cap. It did not recover the full
large-scale HSI method. This within-lineage repetition is consistent with the
structural revisitation documented by \citet{gurkan2026mutation}, but the
present run does not isolate that mechanism from parent selection or prompt
history.

\begin{table}[t]
\centering\small
\begin{tabular}{lrrrr}
\toprule
Executable space & Full cards & Withheld cards & Random & Resample \\
\midrule
Whole module, 24/arm & 0 & 0 & --- & 2 \\
Function patch, 24/arm & 1 & 8 & --- & 0 \\
Typed FBS/DRS, 32/arm & 15 & 4 & 1 & 6 \\
\bottomrule
\end{tabular}
\caption{Earlier ablations: number of fresh candidates solving both small
training cases. The whole-module and patch spaces also included a
no-references arm (1/24 and 0/24 respectively). The typed withheld-card column
is its ordinary-LLM arm. The patch total of eight includes seven variants of
one active-set family; the typed experiment used a different action space and
had a later corrected retirement adapter. Rows are separate experiments, not
pooled hit-rate estimates or demonstrations of complete method recovery.}
\label{tab:llm-loop-earlier-ablation}
\end{table}

A compact typed space removed most source-format failures and let the model
compose textbook FBS or DRS, factor policies, precision, device, retirement,
and work width. In a separate small-case league starting from plain FBS, the
LLM reached DRS with a cached inverse and GPU execution on fresh candidate~6.
Its paired three-repetition contest against the fixed audit ADMM program
showed a $1.157\times$ geometric-mean speed ratio on the \emph{two training
cases}; a paired bootstrap over these fixed cases returned
$[1.141,1.176]$. That interval concerns timing repeatability on those cases,
not population-level generalization. The solver then passed ten of ten held-out
validation executions. It had neither observed pixel retirement nor screened
certificate cadence, so it was a recovery of the mathematical core and factor
policy, not the complete HSI implementation.

The decisive prospective change made certificate policy executable and exposed
a direct ADMM spelling alongside the already available DRS map. For this
identity split, the two updates use different carried variables but have the
same primal iteration under $v=z+u$; a differential test of 300 transitions
across step sizes and relaxations had maximum state discrepancy
$2.14\times10^{-14}$. This equivalence was documented before the full scheduled
method was found. The original direct-ADMM audit is preserved separately from
a secondary semantic audit that accepts this explicit DRS representation.
The direct ADMM adapter is tested numerically against an independent NumPy
implementation, including factor choices and rank-deficient admission, without
claiming a new formal theorem for its execution wrapper.

In the final withheld-card chain, candidate~5 first assembled direct ADMM,
float64 GPU execution, a shared host-computed inverse, relaxation $1.6$,
resolvent step $1$ (equivalently $\rho=1$), and pixel retirement. It certified
all three training cases, but still checked on the host after each chunk.
Candidate~6 retained that splitting, factor reuse, precision, device, and
retirement; it selected screened certification and changed the compiled chunk
from $500$ to $128$. It was the \emph{sixth fresh candidate} in that chain.
The private audit matched its emitted source hash to the trusted compiler output
and observed one factor build, actual compiled ADMM dispatches, GPU array
residence, active-column removal, more device screens than external host checks,
and independent certification of every original column. On Cuprite-4096 it
made 93 ADMM-kernel dispatches, 115 device screens, and 23 host certificates,
with the minimum active batch shrinking to 287 from 4,096. On Cuprite-16384
these counts were 122, 144, and 23, with the active batch shrinking to 824.
These are Python dispatch counts into compiled functions, not physical CUDA
kernel counts or actual iteration counts. The model's prose rationale is never
used as an audit witness.

\begin{table}[t]
\centering\small
\begin{tabular}{clrrl}
\toprule
Fresh try & Main composition & Step/relaxation & Train solved & Missing mechanism \\
\midrule
1 & DRS, float32, host check & $10^{-3}/1.6$ & 1/3 & numerical accuracy, cadence \\
2 & ADMM, float64, host check & $10^{-3}/1.6$ & 1/3 & effective penalty, cadence \\
3 & FBS, float32, host check & $0.5/1$ & 1/3 & splitting, cadence \\
4 & FBS, float32, host check & $0.9/1$ & 1/3 & splitting, cadence \\
5 & ADMM, float64, host check & $1/1.6$ & 3/3 & screened cadence \\
6 & ADMM, float64, device screen & $1/1.6$ & 3/3 & complete recovery \\
\bottomrule
\end{tabular}
\caption{First six distinct withheld-card proposals in order. The first
numeric value is the resolvent step for DRS/ADMM ($1/\rho$) and the forward
step fraction for FBS; the second is relaxation where applicable. All six used
the same legal grammar, fixed three-case training distribution, and an A100
worker. Tries 1, 2, 5, and 6 chose a shared host-computed inverse and active
column retirement; tries 3 and 4 chose FBS. Candidate~6 changed the host
check to a device screen and chunk length from $500$ to $128$ relative to
candidate~5. The six Qwen calls took $22.137$\,s of proposer time, used
$25{,}973$ prompt and $2{,}700$ completion tokens, and produced six fresh
evaluations. The training-case algorithm clocks summed to $1{,}118.485$\,s
across those six candidates, including capped failures; this sum is search
compute, not time-to-solution for the final solver.}
\label{tab:llm-loop-lineage}
\end{table}

Candidates~2 and~5 differed in the canonical genome only at the resolvent
step: $10^{-3}$ versus $1$ (equivalently $\rho=1000$ versus $1$), turning two
over-budget real cases into two certified real cases under the same other
selected moves. Candidates~5 and~6 changed \emph{both} certification policy
and compiled chunk length, so their Cuprite-16384 training-time difference
($47.472$ versus $36.392$\,s) cannot be attributed to cadence alone. These
paired lineage observations explain the search path; they are not a randomized
single-gene ablation. The full-card, no-card, resampling, and random controls
address broader search behavior at matched fresh budgets.

To isolate the two changes in candidate~6, we froze a \emph{post-search}
$2\times2$ factorial on the already inspected Cuprite-10000 sealed batch.
Starting with candidate~5, only the certificate policy (host after every
chunk versus device screen) and compiled chunk length (500 versus 128) changed.
The host-500 and screen-128 source hashes exactly match candidates~5 and~6.
Each combination ran three new $120$\,s-cap processes on the same A100 GPU~2,
in independently shuffled order; all twelve independently certified the full
batch. Because this case had already been used to confirm the search winner,
these numbers diagnose cost and do not count as new transfer evidence or a
second model discovery.

\begin{table}[t]
\centering\small
\begin{tabular}{lrrl}
\toprule
Certificate policy & Chunk & Median (s) & Three full-clock times (s) \\
\midrule
Host every chunk & 500 & 22.244 & 22.244, 21.132, 22.284 \\
Host every chunk & 128 & 29.834 & 29.834, 29.011, 48.839 \\
Device screen & 500 & 18.598 & 18.598, 27.755, 18.335 \\
Device screen & 128 & 19.752 & 20.576, 19.752, 19.730 \\
\bottomrule
\end{tabular}
\caption{Post-search cadence/chunk factorial on one held-out HSI batch.
All twelve runs passed the full original-data certificate. Device screening
reduced the median at both chunk lengths, but screen-500 was slower than
host-500 in one of three paired repetitions. Chunk~128 was not the fastest
screened choice on this batch. These are measured algorithm clocks including
setup and checking, not a cross-distribution performance estimate.}
\label{tab:llm-loop-cadence-factorial}
\end{table}

\begin{table}[t]
\centering\small
\begin{tabular}{lrrr}
\toprule
Input and split & Cases certified & Median algorithm time (s) & Active columns, minimum \\
\midrule
Cuprite-4096, train & 1/1 & 16.606 & 287 \\
Cuprite-16384, train & 1/1 & 36.392 & 824 \\
Synthetic-96, train & 1/1 & 2.513 & 32 \\
Cuprite-6144, sealed & 3/3 & 17.900 & 444 \\
Cuprite-10000, sealed & 3/3 & 22.055 & 772 \\
Synthetic-192, sealed & 3/3 & 3.055 & 128 \\
\bottomrule
\end{tabular}
\caption{Direct ADMM method assembled by the withheld-card LLM loop.
Every sealed run used the independently frozen full-batch float64 certificate;
real held-out pixels do not occur in the search batches. ``Active columns''
comes from observed state, not a declared gene. Algorithm time starts after JAX
runtime initialization but includes solver setup and JIT compilation.}
\label{tab:llm-loop-recovery}
\end{table}

A card-assisted chain also reached a full scheduled composition. Its first
semantic recovery was candidate~3: DRS coordinates, one shared Cholesky factor,
GPU float64, retirement, and screened certification. One Cuprite-4096 training
row initially failed during the SSH banner exchange, before remote execution;
it was explicitly remeasured once, and the original search failure and extra
cost remain in the archive. That recheck, the other two training cases, and
all nine frozen sealed runs passed. Its sealed medians were $40.888$\,s on
Cuprite-6144, $72.061$\,s on Cuprite-10000, and $4.781$\,s on Synthetic-192.
This is a positive control for reachability, not evidence that expert cards
improve the hit rate; the withheld-card chain also recovered the method and was
faster on these particular held-out batches. The completed simple-loop ablation
must use the fixed candidate budgets, including all failures, before any
comparative hit-rate conclusion.

\begin{table}[t]
\centering\small
\resizebox{\linewidth}{!}{%
\begin{tabular}{lrrrr}
\toprule
Arm & Seed-0 fresh & Seed-0 solves / semantic / literal & Seed-1 fresh & Seed-1 solves / semantic / literal \\
\midrule
Full expert cards & 8/8 & 3 / 2 / 2 & 8/8 & 8 / 7 / 1 \\
Withheld expert cards & 8/8 & 3 / 1 / 1 & 8/8 & 2 / 0 / 0 \\
Uniform random & 8/8 & 0 / 0 / 0 & 8/8 & 1 / 1 / 0 \\
Independent model resampling & 8/8 & 0 / 0 / 0 & 8/8 & 1 / 0 / 0 \\
\bottomrule
\end{tabular}%
}
\caption{Completed prospective final-space search on September~26.
``Solves'' counts distinct candidate programs certifying all three training
cases. ``Semantic'' additionally requires the pre-existing observed
ADMM/DRS-equivalent factor/GPU/retirement/screened-cadence audit; ``literal''
requires direct ADMM coordinates. Uniform random seed~1 candidate~7 was a
semantic training hit, so method recovery is not unique to the LLM policy.
Search cases, legal support, fresh budget, and independent
judge are fixed across arms; expert hints are the only extra prompt content in
the full-cards arm.}
\label{tab:llm-loop-final-ablation}
\end{table}

We froze the uniform-random seed~1 candidate~7 before opening sealed data.
Its compiled DRS program passed all nine independent held-out certificates
and the same semantic execution audit, with source hash unchanged. The sealed
median clocks were $27.859$\,s on Cuprite-6144, $31.141$\,s on
Cuprite-10000, and $3.766$\,s on Synthetic-192. This shows that the exposed
finite grammar itself makes a full method reachable by random search in a
small number of trials. The observed LLM and bandit first-hit differences
cannot establish a discovery advantage from two seeds.

The earlier adversarial baseline/champion league asks a different question.
It freezes the manually authored ADMM audit program as an initial opponent,
without showing its source to the proposer. Two sibling proposals share each
prompt. Single executions screen candidates, but only a challenger that
certifies \emph{every} case enters a paired contest against the current
champion, the original fixed anchor, and a rotating older opponent. Contest
order is randomized; each case has three independent repetitions. Promotion
requires the paired bootstrap lower bound from 2,000 resamples to exceed
$1.05\times$ against the champion, no case to regress by more than $20\%$,
and the corresponding tests against the other opponents. The champion then
becomes the next opponent and, after promotion, the next parent. This is
baseline-to-self competition under a fixed certificate; it updates the search
state, not the model weights. In the
small typed league, 22 fresh candidates produced one confirmed promotion,
on candidate~6, and then 16 further candidates produced no promotion. A
separate league continuing from the already discovered CPU active-set solver
had one promotion against the reference after 11 candidates, but that first
promotion did \emph{not} demonstrate an improvement over the inherited
active-set program. Its matched independent-resampling continuation had zero
promotions in eight candidates. On four larger training cases, three league
continuations of twelve fresh candidates each produced zero promotions under
a $60$\,s cap; even the original simple ADMM audit reference missed the
Cuprite-16384 cap. These are negative results for repeated performance
promotion under those budgets, not a contradiction of the later method-recovery
result, which asks whether the model assembled the target mechanism and passes
an independent certificate. The recovery chain used incumbent selection by
solved count and measured time; it did not receive the hidden paper reference
source or a reward for matching a private list of components.

A separate paper-scale replay uses the complete $47{,}750$-pixel Cuprite image
and a $300$\,s algorithm cap. Its original manually authored single-target
schedule reached the certificate in $52.227$, $52.676$, and $51.353$\,s on an
A100. The direct withheld-card recovery certified the same full image in
$56.063$, $55.388$, and $54.974$\,s, all within the cap. Its median was
$55.388$\,s, $1.061\times$ the manual schedule's median on the same A100 slot.
The frozen source and all three execution traces also passed the private
method audit, including actual factor reuse, GPU updates, retirement, and
screened checking.
In its first replay, it made one factor build, 122 compiled ADMM dispatches,
136 device screens, and 15 host certificates; its active batch shrank from
47,750 to 2,274 columns. The card-assisted semantic recovery reached the
numerical certificate in all three replays, with times $303.579$, $299.809$,
and $299.363$\,s; the first is outside the cap and its recorded
\texttt{accepted} flag is therefore false (2/3 within budget). This scale
replay contains training pixels and is a scale check, not another disjoint
holdout. It also differs from the manuscript's original multi-target timing
protocol. The withheld-card method therefore matches the manually assembled
single-target schedule to within $6.1\%$ on this replay, without a performance
superiority claim.

\begin{table}[t]
\centering\small
\begin{tabular}{lrrr}
\toprule
Full-image method & Certificate within 300\,s & Median (s) & Three replay times (s) \\
\midrule
Manual paper schedule & 3/3 & 52.227 & 52.227, 52.676, 51.353 \\
Withheld-card direct recovery & 3/3 & 55.388 & 56.063, 55.388, 54.974 \\
RL-channel semantic recovery & 3/3 & 75.263 & 75.263, 78.001, 65.687 \\
Card-assisted semantic recovery & 2/3 & 299.809 & 303.579, 299.809, 299.363 \\
\bottomrule
\end{tabular}
\caption{Paper-scale replay on the same node-0 A100 GPU~2. The training
distribution already contains pixels from this image, so this is a full-scale
execution check, not a new real-image holdout. All four methods satisfy the
numerical certificate; the 300\,s cap determines the accepted count. Setup,
compilation, transfers, iteration, and certification are included in time.}
\label{tab:llm-loop-fullimage}
\end{table}

\subsection{Next test: learned proposal policy}
The simple loop keeps the proposer weights frozen; its only learning within a
run is the search archive and execution feedback placed in the prompt. Training
an LLM on execution feedback already has precedents in code generation and
algorithm search~\citep{gehring2024rlef,surina2025evolutionrl}; an LLM-weight
update here is therefore a future comparison, not a novelty claim. For a
short-budget, directly executable RL pilot, we selected an online disjoint
LinUCB contextual bandit~\citep{li2010linucb} that controls the \emph{proposal
channel} while keeping Qwen's 27B weights frozen. The three actions are to ask
Qwen for a complete typed design, mutate one legal gene of the current best
design, or sample a fresh legal design uniformly. Mutation has no encoded
target recipe. Every distinct admitted or rejected executable design spends
one fresh evaluation; duplicate and malformed proposals are recorded.

The nine policy features are an intercept, log-normalized design dimensions,
log-normalized right-hand-side count, a GPU indicator, the parent's solved-case
fraction, the measured certificate-to-step cost ratio, an indicator for a
splitting update, and an indicator for screened certification. For each action
$a$ the controller keeps a ridge-initialized matrix $A_a=I$ and vector
$b_a=0$, chooses the action maximizing
$x^\top A_a^{-1}b_a+0.45\sqrt{x^\top A_a^{-1}x}$, then updates
$A_a\leftarrow A_a+xx^\top$ and $b_a\leftarrow b_a+r x$ from the executed
candidate. The fixed bounded reward is at least $0.8$ only when \emph{all}
training cases satisfy the independent numerical judge; accepted designs then
receive a small log-time preference using the full algorithm clock. Otherwise
the reward is at most $0.3$, based on accepted-case fraction and certified
column fraction. The private method-identity audit is never in the reward or
prompt. Model, mutation, and random streams have separate seeds. Source,
requests, outcomes, policy matrices, and selected actions are logged after each
evaluation. This controller is a limited one-step RL policy, not LLM
fine-tuning or evidence of transferable reasoning.

Two on-policy seeds of eight fresh evaluations on the same three HSI training
cases and six-A100 fleet are the direct fixed-budget pilot. The comparison to
the frozen simple loop uses the same compiler, references, numerical gate,
cases, and per-case cap; model-query counts and wall time are also retained.
This controller was chosen after the simple HSI recovery was observed, so the
pilot is an engineering comparison rather than a preregistered RL efficacy
test. A fresh task family is needed for an unbiased transfer claim.
A separate cell-deconvolution assay freezes disjoint LM22 training columns and
LM22/TIL10 held-out columns. It tests an application and distribution shift
under the \emph{same simplex least-squares objective}, not a new mathematical
problem class. The frozen simple loop's second fresh deconvolution proposal
was float64 GPU FBS with screened checking and certified both LM22 training
batches. Freezing that candidate before evaluation, it passed all nine sealed
runs: three each on disjoint LM22-2048, TIL10-2048, and TIL10-4096, with
median times $2.167$, $1.967$, and $1.859$\,s. Its success is an independently
certified solver for this distribution, not recovery of the HSI ADMM method.
The deconvolution RL pilot initializes from a frozen policy after one executed
HSI update; it was launched after the simple deconvolution run had already
found this solver, so its comparison is exploratory and cannot establish an
unbiased transfer gain. Its first two fresh actions (random, then mutation)
failed both training cases. The first model action, on candidate~3, certified
both cases with FBS, and the frozen candidate passed all nine sealed runs.
However, it selected CPU admission, host checks, and a 1,024-step chunk,
whereas the pure-loop candidate used GPU admission and screened checking.
The RL candidate was slower on every held-out batch: median $21.021$ versus
$2.167$\,s on LM22-2048, $3.373$ versus $1.967$\,s on TIL10-2048, and
$5.925$ versus $1.859$\,s on TIL10-4096. Thus it did not improve first-hit
count or certified runtime on this application shift, despite passing the same
numerical bar. At the complete eight-fresh-evaluation budget, six pure-loop
and four bandit candidates certified both deconvolution training batches.
The pure loop made nine model calls (one duplicate proposal), while the bandit
made seven model calls (including duplicate proposals) plus its mutation and
random actions. Fewer model calls did not buy more recoveries here.
The primary scientific quantity remains first complete recovery
within a fixed fresh-evaluation budget on genuinely new classes and hardware.
Neither that transfer nor any RL gain is inferred from a successful HSI-only
run. The H100 serves the proposer; timing on it would require a separate
non-overlapping session after inference completes. All present RL executions
use A100s, so the policy's GPU indicator is constant and these runs cannot
establish learned cross-device adaptation.

HSI bandit seed~0 completed all
eight fresh evaluations: one mutation, one random design, and six model
proposals. The third certified all three training cases, but lacked observed
active-column retirement and therefore failed both the literal and semantic
whole-method audits. The fourth retained screened GPU splitting and factor
reuse, added active-column retirement, certified all three cases, and
passed the pre-existing ADMM/DRS \emph{semantic} execution audit. Its literal
direct-ADMM audit remains false because it executes DRS coordinates.
The controller received rewards $0.100$, $0.100$, and $0.873$ after the first
three executions; the high reward on the third, model-generated solver led it
to query the model again for candidate~4 rather than spend that next fresh
evaluation on random sampling or mutation. Candidate~4 received $0.896$.
This is the observed feedback-to-action path, not an isolated proof that
LinUCB caused the subsequent model proposal to succeed.
This is first training method recovery at candidate~4 versus candidate~6 in
the frozen withheld-card seed~0 chain, a two-evaluation reduction in one
seed. Four of eight bandit candidates passed the pre-existing semantic audit,
versus one of eight in that plain withheld-card chain. A separately seeded
uniform-random chain also produced one semantic hit in eight and passed all
nine sealed checks; this prevents attributing recovery itself to either the
LLM or LinUCB. The bandit made seven
model calls including one duplicate proposal, versus nine model calls
including one duplicate in the plain chain. Up to first recovery, the bandit
used two Qwen calls ($10.964$\,s, $6{,}890$ prompt and $1{,}490$ completion
tokens); the plain chain used six ($22.137$\,s, $25{,}973$ prompt and $2{,}700$
completion tokens). The frozen first RL recovery then
passed all nine sealed checks and all three paper-scale replays; both the
source-hash and semantic execution audits passed. Its sealed medians were
$16.823$\,s on Cuprite-6144, $20.867$\,s on Cuprite-10000, and $3.888$\,s
on Synthetic-192. The full-image median was $75.263$\,s, slower than the
withheld-card direct recovery's $55.388$\,s and the manual schedule's
$52.227$\,s. These are one-seed observations, not a controlled estimate of a
transferable RL effect.
Seed~1 completed eight evaluations (five model, two mutation, one random
design); two candidates certified all three training cases, but none passed
the semantic full-method audit. Thus across the two HSI seeds the bandit had
four method hits in sixteen fresh evaluations, versus one in sixteen for the
withheld-card loop, one in sixteen for uniform random, and zero in sixteen for
independent model resampling. These counts are descriptive: the bandit was
designed after seeing the first plain HSI result, and two seeds do not estimate
a transferable policy benefit. On
deconvolution, the bandit first certified both LM22 training
batches at fresh candidate~3, one candidate later than the frozen Qwen loop's
candidate~2. Its frozen first training solver passed all nine sealed runs but
was slower than the pure-loop solver on every batch, as reported above.
This deconvolution result is negative for the chosen RL controller;
cross-class and cross-device gains remain untested.

At eight fresh evaluations in seed~0, the sums of all measured training-case
algorithm clocks were $1{,}501.056$\,s for the plain withheld-card HSI loop
and $1{,}007.012$\,s for the bandit loop. The analogous deconvolution sums
were $308.896$ and $632.461$\,s. Every candidate in these four chains passed
admission, so these sums contain actual timed training evaluations, including
failed and over-budget runs; they exclude admission, proposer inference,
transport, and lock waiting. HSI model-call time summed to $32.524$\,s for
plain Qwen and $38.680$\,s for bandit-controlled Qwen, despite the latter's
smaller call count. The sums describe search cost on these particular chains,
not the final solver runtime or a cross-task efficiency guarantee.

\clearpage
\begin{table}[t]
\centering\small
\resizebox{\linewidth}{!}{%
\begin{tabular}{lllllll}
\toprule
Loop & Policy weights & HSI first; hits@8 & Sealed HSI & Deconv first; hits@8 & Sealed deconv & New class/device \\
\midrule
Simple typed LLM & frozen Qwen 27B & 6; 1/8 s0; 0/8 s1, semantic & 9/9 & 2; 6/8 & 9/9 & not tested/not tested \\
RL channel policy & LinUCB online; Qwen frozen & 4; 4/8 s0; 0/8 s1, semantic & 9/9; full 3/3 & 3; 4/8 & 9/9, slower & not tested/not tested \\
Uniform random control & no learned weights & 7; 0/8 s0; 1/8 s1, semantic & 9/9 & not run & not run & not tested/not tested \\
\bottomrule
\end{tabular}%
}
\caption{Completed fixed-budget comparison on September~26, 2026. The
seed~0 RL hits are ADMM-equivalent DRS executions; their original literal-ADMM
audits are false. Bandit seed~1 completed eight fresh evaluations with zero
semantic method hits. Uniform random also had one
semantic hit in its completed seed~1 and passed all nine sealed checks.
This one-step RL controller does
not update LLM weights. A transferable RL claim requires fixed-budget recovery
on independent classes and devices, with all proposal and execution costs
reported. Deconvolution is an application shift within the same mathematical
class. ``First'' counts fresh executable candidates; the HSI and deconvolution
success conditions differ as described in the text.}
\label{tab:llm-loop-rl-comparison}
\end{table}
\FloatBarrier
